\SourceLeaf{228}%
Sohn eingeschlagen hat, und die Resultate, zu denen er geführt ist, kommen fast durchgehends mit meinen eigenen, zum Theile schon seit 30—35 Jahren angestellten Meditationen überein. In der That bin ich dadurch auf das Äusserste überrascht. Mein Vorsatz war, von meiner eigenen Arbeit, von der übrigens bis jetzt wenig zu Papier gebracht war, bei meinen Lebzeiten gar nichts bekannt werden zu lassen. Die meisten Menschen haben gar nicht den rechten Sinn für das, worauf es dabei ankommt, und ich habe nur wenige Menschen gefunden, die das, was ich ihnen mittheilte, mit besonderm Interesse aufnahmen. Um das zu können, muss man erst recht lebendig gefühlt haben, was eigentlich fehlt, und darüber sind die meisten Menschen ganz unklar. Dagegen war meine Absicht, mit der Zeit alles so zu Papier zu bringen, dass es wenigstens mit mir dereinst nicht unterginge.

Sehr bin ich also überrascht, dass diese Bemühung mir nun erspart werden kann und höchst erfreulich ist es mir, dass gerade der Sohn meines alten Freundes es ist, der mir auf eine so merkwürdige Art zuvorgekommen ist.

Sehr prägnant und abkürzend finde ich die Bezeichnungen: doch glaube ich, dass es gut sein wird, für manche Hauptbegriffe nicht bloss Zeichen oder Buchstaben, sondern bestimmte Namen festzusetzen, und ich habe bereits vor langer Zeit an Einige solcher Namen gedacht. So lange man die Sache nur in unmittelbarer Anschauung durchdenkt, braucht man keine Namen oder Zeichen; die werden erst nöthig, wenn man sich mit Andern verständigen will. So könnte z.~B. die Fläche, die Dein Sohn $F$ nennt, eine Parasphäre, die Linie $L$ ein Paracykel genannt werden: es ist im Grunde Kugelfläche, oder Kreislinie von unendlichem Radius. Hypercykel könnte der Complexus aller Punkte heissen, die von einer Geraden, mit der sie in Einer Ebene liegen, gleiche Distanz haben; eben so Hypersphäre. Doch das sind alles nur unbedeutende Nebensachen: die Hauptsache ist der Stoff, nicht die Form.

In manchem Theile der Untersuchung habe ich etwas andere Wege eingeschlagen: als ein Specimen füge ich einen rein geometrischen Beweis (in den Hauptzügen) von dem Lehrsatze bei, dass die Differenz der Summe der Winkel eines Dreiecks von $180^{0}$ dem Flächeninhalte des Dreiecks proportional ist.

\Needspace{6\baselineskip}
I. Der Complexus dreier Geraden $ab$, $cd$, $ef$, die so beschaffen sind, dass $ab\TripleParallel dc$, $cd\TripleParallel fe$, $ef\TripleParallel ba$, bildet eine Figur, die ich $T$ nenne. Es lässt
